% 2010 年浙江卷理科第 12 题：三视图（正视图、侧视图、俯视图）。
% 来源 PDF：普通高考/2010/2010浙江理.pdf，SHA-256 c61b26c8cab009d567644272afc34d81e3174ffa90244c333ecfb27d61034a00。
% 原页：第 1 页（A3 三栏）第 2 栏中部，书页码 639。
% 重绘依据：正视图与侧视图均为下底 8、上底 4、高 3 的等腰梯形，其上叠 4×2 的矩形；
%   俯视图为外方 8×8、内方 4×4、四角以对角线相连。尺寸 2、4、2 与 2、3 照原页标注。
\begin{tikzpicture}[scale=0.4, semithick, every node/.style={inner sep=1.5pt, font=\small},
  dim/.style={<->, >=stealth, thin}]

  % ---------- 正视图 ----------
  \draw (-4,0) -- (4,0) -- (2,3) -- (-2,3) -- cycle;
  \draw (-2,3) rectangle (2,5);
  % 顶部尺寸 2,4,2
  \draw[thin] (-4,5.25) -- (-4,5.85);
  \draw[thin] (-2,5.25) -- (-2,5.85);
  \draw[thin] (2,5.25) -- (2,5.85);
  \draw[thin] (4,5.25) -- (4,5.85);
  \draw[dim] (-4,5.55) -- (-2,5.55);
  \draw[dim] (-2,5.55) -- (2,5.55);
  \draw[dim] (2,5.55) -- (4,5.55);
  \node[above] at (-3,5.55) {$2$};
  \node[above] at (0,5.55) {$4$};
  \node[above] at (3,5.55) {$2$};
  % 右侧尺寸 2,3
  \draw[thin] (4.15,5) -- (4.75,5);
  \draw[thin] (4.15,3) -- (4.75,3);
  \draw[thin] (4.15,0) -- (4.75,0);
  \draw[dim] (4.45,3) -- (4.45,5);
  \draw[dim] (4.45,0) -- (4.45,3);
  \node[right] at (4.45,4) {$2$};
  \node[right] at (4.45,1.5) {$3$};
  \node[below] at (0,-0.9) {正视图};

  % ---------- 侧视图 ----------
  \begin{scope}[shift={(11,0)}]
    \draw (-4,0) -- (4,0) -- (2,3) -- (-2,3) -- cycle;
    \draw (-2,3) rectangle (2,5);
    \node[below] at (0,-0.9) {侧视图};
  \end{scope}

  % ---------- 俯视图 ----------
  \begin{scope}[shift={(0,-9)}]
    \draw (-4,0) rectangle (4,-8);
    \draw (-2,-2) rectangle (2,-6);
    \draw (-4,0) -- (-2,-2);
    \draw (4,0) -- (2,-2);
    \draw (4,-8) -- (2,-6);
    \draw (-4,-8) -- (-2,-6);
    % 右侧尺寸 2,4,2
    \draw[thin] (4.15,0) -- (4.75,0);
    \draw[thin] (4.15,-2) -- (4.75,-2);
    \draw[thin] (4.15,-6) -- (4.75,-6);
    \draw[thin] (4.15,-8) -- (4.75,-8);
    \draw[dim] (4.45,-2) -- (4.45,0);
    \draw[dim] (4.45,-6) -- (4.45,-2);
    \draw[dim] (4.45,-8) -- (4.45,-6);
    \node[right] at (4.45,-1) {$2$};
    \node[right] at (4.45,-4) {$4$};
    \node[right] at (4.45,-7) {$2$};
    \node[below] at (0,-8.9) {俯视图};
  \end{scope}
\end{tikzpicture}
